\subsection{Finite definite-theta operations cannot produce the identified input}

\begin{proposition}[P-F68: a finite-operation obstruction and its quantitative cost]
\label{prop:finitethetagrowth}
For a positive-definite quadratic form $Q$ on a real vector space,
a full lattice $L$ in that space, a coset $a+L$, and a polynomial
$p$, call
\[
       \theta_{L+a,p}(z)=\sum_{x\in L+a}p(x)z^{Q(x)}
                                                        \tag{P-F68a}
\]
a weighted definite theta series.  Assume its exponents are rational
with a common denominator.  Finite linear combinations of these series,
multiplied by fixed rational powers of $z$, have polynomially bounded
Fourier coefficients.  This class is closed under finite sums,
products, positive-integer Adams substitutions $z\mapsto z^r$, and
the coefficient derivation $z\,d/dz$.  It is consequently also closed
under any fixed finite composition of the ordinary symmetric- and
exterior-power operations, interpreted through their usual Newton
identities over a characteristic-zero coefficient ring.

No series in this class equals any component of the actual input
$g=\Theta_{E_6}/\delta$ from P-F46, even after restriction to a fixed
positive affine arithmetic progression of its coefficient sequence.
Thus no fixed finite sequence of the indicated operations on definite
theta series supplies that unchanged input.

Here is an explicit bound retaining the data costs.  Suppose a component
is expressed as
\[
              h(z)=\sum_{j=1}^J a_j z^{\beta_j}
                       \theta_{L_j+a'_j,p_j}(z),
                                                        \tag{P-F68b}
\]
where $\dim L_j\le d$, $|\beta_j|\le B$, the least distance between
distinct lattice points in the norm $\|x\|_j=\sqrt{Q_j(x)}$ is at
least $\eta>0$, and
\[
              |p_j(x)|\le C_j(1+\|x\|_j)^k.
\]
Put $A=\sum_j|a_j|C_j$.  For $0\le\nu<1$, $N\ge0$, and
$R=N+B+1$, one has
\[
 |[z^{N+\nu}]h|
 \le A(1+\sqrt R)^k(1+2\sqrt R/\eta)^d.                \tag{P-F68c}
\]
If this coefficient equals $c_\mu(N)=[z^{N+\nu_\mu}]g_\mu$, then
$A>0$ and
\[
 \log_2 A+k\log_2(1+\sqrt R)
       +d\log_2(1+2\sqrt R/\eta)
              \ge\lfloor\sqrt{N+1}\rfloor-1.           \tag{P-F68d}
\]
In particular, for a family of such formulas with
$A\le N^{C_1}$, $B+1\le N^{C_2}$, and
$\eta^{-1}\le N^{C_3}$, where the constants are fixed and $N\ge2$,
matching this one coefficient forces
\[
                     d+k=\Omega(\sqrt N/\log N).
                                                        \tag{P-F68e}
\]
The assertion permits signed or complex summands; it does not rely on
positivity of their coefficients.
\end{proposition}

\begin{proof}
First count lattice points using only Euclidean volume.  In the norm
$\sqrt Q$, distinct points of a coset have distance at least $\eta$.
The open balls of radius $\eta/2$ about those with $Q(x)\le R$ are
disjoint.  They lie in the ball of radius $\sqrt R+\eta/2$ about
zero.  Comparing volumes in dimension $d_j$ gives
\[
 \#\{x\in L+a:Q(x)\le R\}
              \le (1+2\sqrt R/\eta)^{d_j}.             \tag{P-F68f}
\]
The equality of volumes up to the $d_j$th power follows after a
linear change of coordinates taking $Q$ to the standard squared
norm.  Using open rather than closed balls eliminates boundary
intersections; boundaries have zero volume.  A rank-zero lattice
has at most one point, giving the same bound with exponent zero.

For a term of (P-F68b) contributing at exponent $N+\nu$, its vectors
satisfy $Q_j(x)=N+\nu-\beta_j\le R$.  If that number is negative
there is no contribution.  Otherwise every summand has absolute
value at most $C_j(1+\sqrt R)^k$.  Multiply by (P-F68f), sum with
the triangle inequality, and use $d_j\le d$ to obtain (P-F68c).
Every polynomial admits the stated bound for some $C_j$ and its
degree $k_j$, since all norms on its finite-dimensional space are
equivalent.  The lattice has a positive least distance.  Thus
fixed finite data give polynomial growth in $N$.

The closure assertions are literal identities.  A product of two
weighted theta series is the theta series of the direct-sum lattice,
with quadratic form $Q_1+Q_2$ and polynomial $p_1p_2$; its coefficients
are finite sums because the forms are positive definite.  The Adams
substitution replaces $Q$ by $rQ$ and multiplies its exterior exponent
shift by $r$.  Differentiating a shifted term gives
\[
 z\frac d{dz}\left(z^\beta\theta_{L+a,p}(z)\right)
       =z^\beta\theta_{L+a,(\beta+Q)p}(z),              \tag{P-F68g}
\]
so it only changes the polynomial weight.  Finite sums and scalar
multiples are already allowed.  Finally the character identities
\[
 \sum_{r\ge0}\sigma^r(f)t^r
   =\exp\left(\sum_{m\ge1}\psi^m(f)t^m/m\right),\qquad
 \sum_{r\ge0}\lambda^r(f)t^r
   =\exp\left(\sum_{m\ge1}(-1)^{m-1}\psi^m(f)t^m/m\right)
                                                        \tag{P-F68h}
\]
show that each fixed coefficient $r$ is a polynomial over $\mathbb Q$
in $\psi^1(f),\ldots,\psi^r(f)$.  These are the ordinary Newton
identities, with $\psi^m(z)=z^m$, already used in the plethystic
bridge.  They therefore use only the proved finite closure operations.
This statement is about a fixed $r$ or a fixed finite composition,
not the entire infinite generating series in (P-F68h).

For the specified input, the elementary subset-of-parts injection
(P-F46i)--(P-F46j) gives, for every component and $N\ge0$,
\[
                  c_\mu(N)\ge 2^{\lfloor\sqrt{N+1}\rfloor-1}.
                                                        \tag{P-F68i}
\]
Taking logarithms of (P-F68c) and (P-F68i) proves (P-F68d).
For fixed data, the right side of the logarithm of (P-F68c) is
$O(\log(N+2))$, whereas the right side of (P-F68d) has order
$\sqrt N$.  Equality for all sufficiently large $N$ is impossible.
On any fixed positive affine progression $N=as+b$, the same
comparison is between $O(\log(s+2))$ and a positive multiple of
$\sqrt s$; clearing a fixed denominator covers rational affine
reindexings on their compatible progressions.

For the quantitative assertion, the stated polynomial bounds imply
$R\le2N^{\max\{1,C_2\}}$ and, with constants depending only on
$C_1,C_2,C_3$,
\[
 \log_2 A\le C_1\log_2N,\qquad
 \log_2(1+\sqrt R),\ \log_2(1+2\sqrt R/\eta)
                                            \le C_4\log_2N.
\]
If necessary replace $C_2,C_3$ by nonnegative larger constants;
this weakens their hypotheses and makes the displayed bounds valid
uniformly for $N\ge2$.  Equation (P-F68d) now gives
$C_1+C_4(d+k)\ge(\lfloor\sqrt{N+1}\rfloor-1)/\log_2N$,
which proves (P-F68e).
\end{proof}

\paragraph{The precise implication for ordinary compact-source theta chains.}
The theorem applies to an ordinary final theta integral to the
rank-one symplectic group, or its specified $U(1,1)$ restriction
as in P-F51, from a fixed positive-definite rational unitary group
with fixed finite level and
finite Gaussian-polynomial Schwartz data, after its usual holomorphic
archimedean normalization.  This is the class of explicit data used
for the positive factor in P-F51.  Here is the reduction, to make its
scope testable.  The compactness assertion is an application of
Borel, \emph{Some finiteness properties of adele groups over number
fields}, Theorem~5.6(ii), printed p.~21:
\url{https://www.numdam.org/item/PMIHES_1963__16__5_0.pdf}.
Its hypotheses hold here.  A unitary matrix for a positive-definite
form is diagonalizable over $\mathbb C$, so its only unipotent
element is the identity.  A rational character of the unitary group
maps its connected compact real group into $\mathbb R^\times$,
hence has image $\{1\}$; the real group is Zariski dense, so the
character is trivial.  Inversion changes Borel's right quotient
to the left quotient used here.  Modulo the real compact group and
an open finite-level subgroup, this compact quotient has a discrete,
therefore finite, quotient.  Borel's Theorem~5.1, printed p.~19,
also gives finite double classes directly; an arbitrary compact
open level is commensurable with the integral level there.
A finite-adelic
Schwartz function is a finite linear combination of lattice-coset
indicators after restriction to rational vectors: local constancy
supplies a common open additive lattice, and compact support permits
only finitely many of its cosets.  On each source double class the
finite-adelic transformation changes those lattices and cosets, giving
still finitely many such terms.

At the real place a finite Gaussian-polynomial vector has a polynomial
of fixed bounded degree times the invariant Gaussian.  Compact-unitary
translation preserves the Gaussian and that finite-dimensional
polynomial space.  Integrating it against an integrable fixed source
function therefore replaces its finitely many polynomial coefficients
by scalar integrals; it does not create an infinite-dimensional space
of polynomial weights.  More explicitly, decompose the polynomial into its finitely many
homogeneous pieces of degrees $\ell$.  The Schr\"odinger action of
$m(\sqrt y)$ contributes $y^{d_j/4+\ell/2}$ on such a piece,
while the Gaussian and unipotent action give
$e^{2\pi i\tau Q_j(x)}$.  Removing that specified factor on each
weight component leaves precisely (P-F68a).  This is the
normalization intended here; no arbitrary Schwartz vector is
being declared to have this finite polynomial expansion.  Absolute convergence
justifies the interchange: on each compact source class, polynomial
growth of the weights is uniformly dominated by the Gaussian.
Thus the final holomorphic coefficient sequence is of form (P-F68b).
The real actions used here are explicitly
$(\omega(m(a))\Phi)(x)=a^{d_j/2}\Phi(ax)$ for $a>0$ and
$(\omega(n(u))\Phi)(x)=e^{2\pi iuQ_j(x)}\Phi(x)$,
exactly as fixed in P-F51.  Their displayed homogeneous-polynomial
consequence is a direct substitution.  For the finite-place
Schr\"odinger conventions and orthogonal splitting we use Kudla's
\emph{Notes on the local theta correspondence},
Definition~II.4.2 and Proposition~II.4.3, pp.~37--38, and the
orthogonal splitting on p.~39:
\url{https://www.math.toronto.edu/~skudla/castle.pdf}.

Any fixed finite earlier sequence of ordinary compact-source lifts
which supplies such a final source function and retains these fixed
level/vector hypotheses is subject to the same conclusion.  Its
intermediate representation labels do not change the final lattice
count.  Consequently this concrete finite-theta input dictionary
cannot produce $\Theta_{E_6}/\delta$.  It must introduce an operation
outside this class, such as the identified infinite symmetric
completion for $\delta^{-1}$, or use data which vary without the
fixed bounds.  If a proposed varying chain has
$d+k\le C_0+C_5m$ after $m$ steps and satisfies the polynomial data
bounds in (P-F68e), matching the coefficient at $N$ requires
$m=\Omega(\sqrt N/\log N)$.  This is a conditional resource bound
for that precise construction, not an estimate for the supplied
number-theoretic proof, whose data have not been identified with it.

The proposition neither excludes directly producing the discriminant
output by a different theta construction nor identifies the known
regularized multiplicative lift with the supplied iterated lifts.
It does not cover indefinite or regularized theta kernels, an infinite
sum of levels or ranks, or arbitrary archimedean Schwartz vectors.
The completed symmetric and exterior constructions in P-F53 remain
valid.  The remaining actual multiplicity and source-operation
comparison is stated separately in O-THETA.
